This appendix supplies the constructions used in (3.1)--(3.3).
Throughout, $E=E_\ell$ is fixed, with the lattice, parametrization and
period convention of Section~\ref{sec:rc}. Constants may depend on these
data and on the prescribed finite list of local types. They do not depend
on the number or sizes of the later good conductor primes.

\subsection{One retained class, including discriminant \texorpdfstring{$4h$}{4h}}
All numbered primary references in this subsection are to Kato~\cite{kato}.

Fix a non-CM elliptic curve $E/\mathbf Q$ of conductor $N$, its weight-two newform $f$,
$T=T_2E$, and $V=T\otimes\mathbf Q_2$. In the application
$E=E_\ell$ is fixed first. Put
$\Lambda=\mathbf Z_2[[\Gamma]]=\mathbf Z_2[[t]]$, where $\Gamma$
is the real cyclotomic Galois group. Let
$\kappa:G_{\mathbf Q}\to\Gamma\subset\Lambda^\times$ be the quotient map,
and set $A=\Lambda[1/2]$.
We use arithmetic Frobenius throughout. The action on the cyclotomic
group-ring coefficient is the inverse deck action.

\paragraph{The precise output.}
\begin{proposition}[One retained Kato system and its two comparisons]
\label{ex:app-kato-system}
Use the fixed curve, lattice and cyclotomic convention just specified.
Fix $e\ge\max\{3,v_2(N)\}$ large enough for the fixed initial local data;
the smoothing choice is (K5).
Let $G$ be a finite elementary abelian two-group, let
$\psi:G_{\mathbf Q}\to G$ be a continuous homomorphism, and let
$\chi_x=\chi_0(x\circ\psi)$, for $x\in\widehat G$, be a finite
family of quadratic characters with one fixed real sign and fixed
local squareclasses at the primes dividing $2N$.
Since the characters of $G$ separate
its elements, $\psi(c)=1$ for complex conjugation $c$, and
$\psi|_{G_{\mathbf Q_v}}=1$ at every $v\mid2N$. All varying conductor primes
are good odd primes for $E$. Let $F$ be the squarefree least common
multiple of the odd parts of the primitive conductors. There is a
single retained-smoothed class
\[
 Z_F\in 2^{-c_E}H^1\bigl(G_{\mathbf Q,S},
        T\otimes\Lambda[G](\chi_0\psi\kappa^{-1})\bigr),
 \qquad S=\operatorname{prime}(2NF),                         \tag{K1}
